\section{Introduction}
\label{sec:introduction}
An aerial robot and a ground mobile manipulator offer complementary capabilities:
the former can change its observation position without following a ground route,
whereas the latter can support a manipulator and physically retrieve an object.
Exploiting this complementarity requires more than detecting a target from the
air. The ground robot must eventually occupy a suitable base pose, see the
target at close range, and execute a collision-aware grasp and lift. Information
that improves a map is therefore not necessarily information that makes the
retrieval task executable.

Next-best-view (NBV) planning provides a natural mechanism for deciding where
to observe. Exploration and reconstruction approaches select measurements that
reveal unmapped space or incompletely observed surfaces
\cite{bircher2016nbv,border2018see}. Such objectives address their intended
mapping tasks, but do not by themselves express the capabilities of a different
robot that will act later. For retrieval, a large poorly observed area outside
useful manipulator base placements can contribute substantial generic observation
gain while providing little evidence about an executable manipulation pose.
Conversely, a smaller unresolved region intersecting a promising base footprint
may prevent the entire task from proceeding. With a finite sensing budget, this
distinction makes the spatial allocation of observations consequential.

Task-oriented view selection is not new: prior work selects views for grasp
detection or decides when to stop observing and grasp
\cite{gualtieri2017viewpoint,breyer2022target}. We address a different coupling.
The aerial sensing robot does not execute the grasp, and the region that needs
observation is not restricted to the object's visible surface. It depends on
where the ground manipulator can stand, the area occupied by its base, and the
operational constraints between that stance and a completed retrieval. The
research question is whether these downstream manipulation requirements can
guide aerial observation more effectively than a shared, information-centric
NBV objective under the same observation budget.

Our central idea is to propagate manipulation support backward into the sensing
objective. Given an aerially perceived grasp target, an inverse-reachability
query produces candidate base poses. Validated candidates supply a
joint-margin-based manipulation relevance field. Their exact base footprints
project this relevance into environment space, where it weights the current
observation deficit. A finite-scan acquisition model then estimates which cells
a candidate aerial view can observe during the next accepted window. The policy
selects a view using the resulting task-weighted gain and a shared motion-cost
proxy. After each observation, the system rechecks whether an exact base
candidate is sufficiently observed and passes a common execution screen. This
connects observation choice to a physical handoff rather than treating a heatmap
or a reachable grasp as the final outcome.

We make three contributions:
\begin{enumerate}
\item A manipulation-aware active-perception framework for heterogeneous
aerial-ground retrieval, in which the ground robot's future operation informs
the aerial robot's sensing decisions.
\item A footprint-aware task relevance formulation that combines validated
manipulator reachability, exact base-pose support and operational geometry with
environmental observation deficit, while keeping unsupported capability,
unobserved ground and observed collision evidence distinct.
\item Closed-loop physics-based validation on P450, BUNKER, AUBO i5 and AG95,
using 96 independent scenes and 192 primary paired trials with physical retrieval
success as the primary endpoint. A controlled Generic/Ours comparison shares
the sensing model, candidates, costs, budget and execution pipeline, isolating
the observation-gain weighting.
\end{enumerate}
The proposed contribution is neither a new inverse-kinematics solver nor the
shared platform corrections. The evaluation supports a success-rate advantage
within the defined simulated task population; efficiency, failure mechanisms
and transfer limitations are reported separately.
